\makeatletter
\def\input@path{{./}{../}{../../}{../../../}{../../../../}{preamble/}{../preamble/}{../../preamble/}{../../../preamble/}}
\makeatother
\input{preamble_exams.tex}
\input{graphs/highlights.tex}
% \GTenHighlightsfalse % Uncomment to disable instructional highlights.
\input{graphs/specific_highlights.tex}

\hypersetup{pdftitle={Grade 10 Global Economics: Term 3, Lesson 1, C3},pdfsubject={Home practice: Probability Rules across Models}}
\title{Grade 10 Global Economics\\Term 3, Lesson 1: Home Practice\\C3 --- Probability Rules across Models}
\author{}
\date{}
\begin{document}
\maketitle
\setcounter{tocdepth}{1}
\setcounter{secnumdepth}{2}
\phantomsection\label{toc}
\tableofcontents
\section*{Evaluated criteria}
\begin{description}
\item[\textbf{C3}] Applies probability rules for simple, combined, and complementary events across discrete and continuous probability models to calculate and interpret likelihoods in economic and financial contexts.
\end{description}
\textbf{How to practise.} Translate words into event notation. Before adding or multiplying, state whether events overlap and whether independence is given. Try each task on paper before opening its solution link. Record the model, working, units and an interpretation; use the arrows to return.
All datasets and economic or financial models in this activity are illustrative, rather than observed market evidence.

\clearpage
\section{Introduction}\label{sec:introduction}
\subsection*{Highlight legend}
\phantomsection\label{tab:highlight-legend}
\begin{tabularx}{\linewidth}{@{}>{\raggedright\arraybackslash}p{.20\linewidth}>{\raggedright\arraybackslash}p{.22\linewidth}>{\raggedright\arraybackslash}X@{}}
\toprule
\textbf{Symbol or notation} & \textbf{Concept} & \textbf{Brief meaning}\\
\midrule
$\HLMathEvent{A\cap B}$ & \HLMathEvent{Intersection} & Both events occur. \\
$\HLMathEvent{A\cup B}$ & \HLMathEvent{Union} & At least one occurs. \\
$\HLComplementProbability{P(A^c)}$ & \HLComplementProbability{Complement} & Probability of not A. \\
$P(B\mid A)$ & Conditional probability & Probability after A is known. \\
$P(A\cap B)=P(A)P(B)$ & Independence & Knowing one does not change the other. \\
\bottomrule
\end{tabularx}
\subsection{Introductory example 1: Two customer attributes}\label{ex:intro-one}
For a randomly selected customer, let $\HLMathEvent{A}$ mean pays by card and $\HLMathEvent{B}$ mean uses a voucher. The \HLMathEvent{intersection} $A\cap B$ means both; the \HLMathEvent{union} $A\cup B$ means at least one. A card-paying voucher user belongs to both, so adding $P(A)$ and $P(B)$ alone would count that customer twice.
\subsection{Introductory example 2: A continuous delivery window}\label{ex:intro-two}
Let $\GTenHLValue{T}\sim U(0,10)$ hours be an illustrative delivery delay. For $D=\{T>8\}$, $\HLProbability{P(D)=0.2}$ and the \HLComplementProbability{complement probability} is $\HLComplementProbability{P(D^c)=0.8}$. The complementary event is delivery within eight hours. Exact endpoints have zero probability in this model.
\section{Notation Reference and Practice Procedure}\label{sec:notation-and-procedure}
\subsection{Notation reference}\label{tab:notation-reference}
\begin{tabularx}{\linewidth}{@{}>{\raggedright\arraybackslash}p{.16\linewidth}>{\raggedright\arraybackslash}p{.18\linewidth}>{\raggedright\arraybackslash}p{.27\linewidth}>{\raggedright\arraybackslash}X@{}}
\toprule
\textbf{Symbol or notation} & \textbf{Concept} & \textbf{Definition} & \textbf{Meaning in the introductory examples}\\
\midrule
$\HLMathEvent{A\cap B}$ & \HLMathEvent{Intersection} & Outcomes common to two events. & \textbullet\ \hyperref[ex:intro-one]{Example 1}: Card payment together with voucher use. \\
$\HLMathEvent{A\cup B}$ & \HLMathEvent{Union} & Outcomes in either event, including their overlap. & \textbullet\ \hyperref[ex:intro-one]{Example 1}: Card payment or voucher use or both. \\
$\HLComplementProbability{P(A^c)}$ & \HLComplementProbability{Complement} & $P(A^c)=1-P(A)$. & \textbullet\ \hyperref[ex:intro-two]{Example 2}: $D^c$ is completion within eight hours. \\
$P(B\mid A)$ & Conditional probability & For $P(A)>0$, $P(Acap B)/P(A)$. & \textbullet\ \hyperref[ex:intro-one]{Example 1}: Voucher-use probability among card payers. \\
$P(A\cap B)=P(A)P(B)$ & Independence & Equivalent to $P(Bmid A)=P(B)$ when defined. & \textbullet\ \hyperref[ex:intro-one]{Example 1}: Cannot be assumed just because customer attributes differ. \\
\bottomrule
\end{tabularx}
\subsection{Practice procedure}
\begin{enumerate}
\item Define the experiment and events; identify whether probabilities come from discrete masses, frequencies or continuous areas.
\item Use $P(A^c)=1-P(A)$ and $P(A\cup B)=P(A)+P(B)-P(A\cap B)$. Add directly only for disjoint events.
\item Use $P(A\cap B)=P(A)P(B\mid A)$. Replace the conditional probability by $P(B)$ only when independence is justified.
\item Distinguish mutually exclusive events, which cannot occur together, from independent events. Interpret the result for the original experiment.
\end{enumerate}
\section{General Worked Problems}\label{sec:general-worked}
\SectionProblemsTableOfContents{
\SectionProblemLink{prob:sec:general-worked0}{Late deliveries: discrete masses}\\
\SectionProblemLink{prob:sec:general-worked1}{Overlapping intervals in a uniform model}\\
\SectionProblemLink{prob:sec:general-worked2}{Card and voucher decisions}}

\subsection{Late deliveries: discrete masses}\label{prob:sec:general-worked0}
The number $W$ of late shop deliveries in a day has distribution
\[
\begin{array}{c|rrrr}w&0&1&2&3\\ \hline P(W=w)&0.1&b&2b&0.3\end{array}.
\]
Let $A=\{W\ge2\}$ and $B=\{W\text{ is odd}\}$.
\QuestionsTasks[sec:general-worked]{{3}/{sol:sec:general-worked0}/{Find the missing parameter and calculate the probabilities of A, their intersection, their union, and the complement of A.}}
\subsubsection*{Worked solution}\label{sol:sec:general-worked0}
Discrete probabilities must sum to one: $0.1+b+2b+0.3=1$, so $3b=0.6$ and $b=0.2$. Consequently
\[
P(A)=0.4+0.3=\HLProbability{0.7},\quad P(B)=0.2+0.3=0.5.
\]
The common outcome is $W=3$, so $P(A\cap B)=0.3$. The union rule gives
\[
P(A\cup B)=0.7+0.5-0.3=\HLProbability{0.9},\qquad
P(A^c)=1-0.7=\HLComplementProbability{0.3}.
\]
Thus $90\%$ of modeled days have either at least two late deliveries or an odd count, and $30\%$ have fewer than two. Since $0.3\ne0.7(0.5)$, these events are not independent.
\secInicio[prob:sec:general-worked0][sec:general-worked]


\subsection{Overlapping intervals in a uniform model}\label{prob:sec:general-worked1}
Let delivery delay $T\sim U(0,10)$ hours. Define $A=\{2\le T\le7\}$ and $B=\{5\le T\le9\}$.
\QuestionsTasks[sec:general-worked]{{3}/{sol:sec:general-worked1}/{Calculate the intersection, union and complement-of-union probabilities. Explain why direct addition fails and decide whether multiplying the event probabilities is justified.}}
\subsubsection*{Worked solution}\label{sol:sec:general-worked1}
The constant density is $1/10$ per hour. The overlap is $[5,7]$, with area $2/10=0.2$. Hence
\[
P(A)=5/10=0.5,\quad P(B)=4/10=0.4,
\]
\[
P(A\cup B)=0.5+0.4-0.2=\HLProbability{0.7},\quad
P((A\cup B)^c)=1-0.7=\HLComplementProbability{0.3}.
\]
The union is $[2,9]$ and its complement within the support is below two or above nine hours. Direct addition double counts the overlap. Multiplication $0.5(0.4)=0.2$ happens to equal the intersection here, so these particular events \emph{are} independent. That equality must be established; overlap alone neither proves nor disproves independence. Direct multiplication is justified after this calculation, rather than as an initial assumption.
\secInicio[prob:sec:general-worked1][sec:general-worked]


\subsection{Card and voucher decisions}\label{prob:sec:general-worked2}
For one randomly chosen customer, $P(A)=0.60$ for card payment, $P(B)=0.40$ for voucher use, and $P(B\mid A)=0.50$. A different model explicitly assumes that failures of two payment terminals are independent, each with probability $0.10$.
\QuestionsTasks[sec:general-worked]{{3}/{sol:sec:general-worked2}/{Find both-customer-attribute and either-attribute probabilities; decide whether the attributes are independent. Then find the probability that at least one terminal fails.}}
\subsubsection*{Worked solution}\label{sol:sec:general-worked2}
For the customer attributes use the conditional multiplication rule:
\[
P(A\cap B)=0.60(0.50)=\HLProbability{0.30},
\qquad P(A\cup B)=0.60+0.40-0.30=0.70.
\]
They are dependent since $P(B\mid A)=0.50\ne0.40=P(B)$. They are not mutually exclusive because their intersection is positive.
For the terminals the stated independence allows multiplication of the no-failure probabilities:
\[
P(\text{at least one fails})=1-(0.90)(0.90)=\HLProbability{0.19}.
\]
Thus the modeled chance of one or more terminal failures is $19\%$. Two positive-probability mutually exclusive events cannot be independent: their intersection is zero but the product of their probabilities is positive.
\secInicio[prob:sec:general-worked2][sec:general-worked]
\section{Advanced Problems}\label{sec:advanced}
\SectionProblemsTableOfContents{
\SectionProblemLink{prob:sec:advanced0}{Linear density and a changed dependence assumption}}

\subsection{Linear density and a changed dependence assumption}\label{prob:sec:advanced0}
A shipment surcharge $X$ in dollars has density $f(x)=x/18$ on $[0,6]$, zero elsewhere. Let $A=\{X<3\}$ and $B=\{X>4\}$. Separately, a payment has failure probability $0.10$. Initially assume payment failure is independent of $A$; a revised model states $P(\text{failure}\mid A)=0.20$.
\QuestionsTasks[sec:advanced]{{3}/{sol:sec:advanced0}/{Find the probability of A or B and of neither event. Compare the chance of both A and payment failure under the two payment assumptions.}}
\subsubsection*{Worked solution}\label{sol:sec:advanced0}
Find areas geometrically. The lower region is a triangle; the upper region is a trapezium:
\[
P(A)=\tfrac12(3)(3/18)=1/4,\qquad
P(B)=2(4/18+6/18)/2=5/9.
\]
The intervals are disjoint, so $P(A\cap B)=0$ and
\[
P(A\cup B)=1/4+5/9=29/36,\quad
P((A\cup B)^c)=1-29/36=7/36.
\]
About $80.6\%$ of surcharges are below three or above four dollars. A and B are mutually exclusive and are not independent because both have positive probability.
For payment failure, the initial model gives $P(A\cap\text{failure})=(1/4)(0.10)=0.025$. The revised conditional model instead gives $(1/4)(0.20)=0.050$. Dependence doubles this joint probability; the surcharge's own density remains the same.
\secInicio[prob:sec:advanced0][sec:advanced]

\end{document}

